\documentclass[a4paper]{article}

\usepackage[english]{babel}
%\usepackage[utf8x]{inputenc}
\usepackage{amsmath}
\usepackage{graphicx}
\usepackage[colorinlistoftodos]{todonotes}
\usepackage[all]{xy}
\usepackage{dsfont}
\usepackage{amssymb}
\usepackage{bbold}
\usepackage{amsthm}

\usepackage{tikz}
\usepackage{verbatim}
\usetikzlibrary{matrix}
\usetikzlibrary{cd}

\usepackage[autostyle=true]{csquotes}
\usepackage{pdfpages}
\usepackage{hyperref}
\usepackage{fancyhdr}



\newtheorem{theorem}{Theorem}[section]
\newtheorem{proposition}{Proposition}[theorem]
\newtheorem{lemma}{Lemma}[theorem]
\newtheorem{corollary}{Corollary}[theorem]
\newtheorem*{definition}{Definition}


\newcommand{\norm}[1]{\left\lVert#1\right\rVert}

\newcommand{\Rn}{\mathbb{R}^n}


\newcommand{\too}[1]{\overset{#1}{\longrightarrow}}
\newcommand{\hookar}[1]{\overset{#1}{\hookrightarrow}}
\newcommand{\colim}[1]{ \underset{#1}{\texttt{colim}} \text{ } }
\newcommand{\fset}[1]{ \{ #1_0 , \dots , #1_n \} }
\newcommand{\ddt}{\frac{d}{dt}}


\newcommand{\N}{\mathbb{N}}
\newcommand{\Z}{\mathbb{Z}}







\fancyhead[L]{\includegraphics[scale=0.1]{UFT.png}}

\begin{document}


\thispagestyle{fancy}
\vspace*{1cm}

\begin{minipage}[r]{12cm}
    \begin{flushright}
\begin{tikzpicture}
\draw[green, very thick] (0,0) rectangle (3,2);
\end{tikzpicture}
    \end{flushright}
\end{minipage}

\begin{center}
\Large{\textbf{Control 1\\ Álgebra Lineal\\ }}
\end{center}
\vspace{1cm}
\begin{minipage}[r]{15cm}
\begin{flushleft}

Fecha: 5 de octubre 2023 \\%
Duración: 120 minutos  \\
Profesores: Camila Muñoz.\\\hspace{1.7cm} Adrián Silva.\\
\end{flushleft}
\end{minipage}
\vspace{2cm}
\begin{center}
\Large{\textbf{Facultad de Ingeniería}}
\end{center}
\vspace{1cm}
\textbf{En caso de que un alumno sea sorprendido en actitudes deshonestas, ya sea intentando o
dejándose copiar, será calificado con nota mínima 1.0 y perderá toda posibilidad de eliminar
dicha nota. Además, esta conducta quedará archivada como antecedente en su carpeta y
podrá ser causal de eliminación.
Así mismo, está prohibido salir de la sala, usar o mantener sobre la mesa teléfonos celulares
o cualquier tipo de aparato electrónico, a no ser que el profesor lo requiera.}\\



\vspace{5cm}
\makebox[12cm]{\textbf{Nombre Alumno}:\ \hrulefill} 
\newpage
\thispagestyle{empty}
\section*{Preguntas}
\begin{enumerate}


\item {\bf (10 puntos)} Si $A=\begin{pmatrix} 1 & -1 \\ 1 & 2 \end{pmatrix}$, $B=\begin{pmatrix} 0 & 1 \\ 3 & -2 \end{pmatrix}$ y $C=\begin{pmatrix} 1 & 4 \\ -1 & 2 \end{pmatrix}$, resuelva la siguiente ecuación matricial para $X$:
    \[AX^tB+I_2=C.\]
    

    
\newpage

\newpage
\item {\bf (12 puntos)} Utilizando propiedades de determinantes, determine el valor de $m$, si se sabe que 

\[\begin{vmatrix} 1 & 1  & 2 & 1 \\ 3 & -1 & 2& m\\ 2 & -2 & 0& -1\\ 1 & 1 & -1 &  1  \end{vmatrix}=4. \]

   \newpage 
       \item  Considere las matrices $A,D,Q \in \mathcal{M}_3(\mathbb{R})$, donde $Q$ es invertible y además,  ́$ A = Q^{-1}DQ$. Demuestre, mediante propiedades de matrices y determinantes que:

\begin{enumerate}
    \item {\bf (3 puntos)} $A^3=Q^{-1}D^3Q$.
    \item {\bf (3 puntos)} $A^n=Q^{-1}D^nQ$ para todo $n\in \mathbb{N}$.
    \item {\bf (5 puntos)} $|A|=|D|$.
\end{enumerate}

\newpage
    \item {\bf (12 puntos)} Calcule la inversa de
\[\begin{pmatrix} 1 & 1  & 0  \\ -1 & 1 & 2\\ 1 & 0 & 1 \end{pmatrix}, \]
por el m\'etodo $A^{-1}=\dfrac{1}{\det(A)}\operatorname{adj}(A)$.
    \newpage
\newpage

    \item {\bf (15 puntos)} Usando el teorema del rango, determine los valores de $\alpha$ de modo que el sistema
    \begin{align*}
x + y +z &= 4 \\
2x + 5y - \alpha z &= 3\\
x + 2y +  z &= 8
\end{align*}

    tenga solución única, infinitas soluciones o ninguna solución. Para $\alpha=3$, determine el conjunto soluci\'on del sistema.
        
        









\end{enumerate}


\end{document}